\documentclass[main.tex]{subfiles}
\begin{document}

In the U.S., Visa and Mastercard (MC) process 80\% of card transactions and earn operating margins above 60\% \parencite{Visa2020}.
These facts have motivated two decades of litigation and legislation aimed at lowering the fees that merchants pay to accept cards, which totaled \$120 billion in 2019 \parencite{Nilson2020}. 
Key policy interventions include the 2011 Durbin Amendment, which capped the interchange fees merchants pay to accept debit cards, the 2022 proposed Credit Card Competition Act \parencite{Andriotis2022}, and the 2024 DOJ lawsuit against Visa \parencite{Au-Yeung2024}.

But because payment markets are two-sided, standard intuitions on the effects of competition and regulation can be reversed.
Networks tax merchants to subsidize consumers with rewards. 
In a two-sided market, capping merchant fees can reduce potentially desirable consumer rewards \parencite{Rochet.Tirole2003}. 
Conversely, competition can incentivize networks to raise merchant fees to fund greater consumer rewards \parencite{Edelman.Wright2015}.
But despite the rich set of theoretical possibilities, empirical evidence remains scarce.

I develop and estimate a quantitative model of platform competition to identify the sources of welfare losses in U.S. payment markets.
The central problem is not the level of competition but its direction.
Many consumers carry cards from only one network, while almost all merchants accept all major cards to avoid losing sales.
This creates a ``competitive bottleneck'' \parencite{Armstrong2006}in which networks compete by inflating rewards instead of cutting fees.
The result is a high-fee, high-reward equilibrium.
Fee caps at the cost of cash correct this distortion and raise welfare by \$30 billion a year.
But more competition makes things worse. 
Because consumers are more reward-sensitive than merchants are fee-sensitive, network competition inflates both sides and reduces welfare by \$20 billion.
Reduced network competition would reduce rewards and the excessive adoption of high-fee payment methods.
The direction of competition matters more than its intensity. 
Dual routing requirements that increase consumer multi-homing force networks to compete for merchants, reducing merchant fees and rewards and increasing welfare by \$10 billion.

Merchants in the U.S. typically charge consumers uniform prices for purchases made with different payment methods \parencite{Stavins2018}.\footnote{
    This occurs even though cash discounts and card surcharges are largely legal. I explore surcharging both theoretically and empirically in Online Appendix \ref{sec:oa-price-coherence}.
}
This pricing rigidity makes the market two-sided. 
When credit card networks tax merchants to subsidize consumers, their consumers benefit from the rewards but only internalize part of the cost of higher retail prices \parencite{Felt.etal2020}.


I present three reduced-form facts that show a competitive bottleneck in payment markets.
First, consumers are sensitive to rewards.
I use a bank-level panel of payment volumes to show that a 25 basis point reduction in debit rewards in 2010 reduced debit card spending by around 30\%.
Second, accepting cards significantly increases merchant sales.
I present event-study evidence that accepting credit cards increases sales to card-carrying consumers by around $12\%$ for merchants on the margin of acceptance.
Third, I show that while most merchants accept all networks, many consumers only carry credit cards from one network.
Merchants risk lost sales when they decline consumers' preferred payment methods.

In the structural model, payment networks compete as two-sided platforms and set both merchant fees and consumer rewards.
Consumers adopt the best bundle of payment methods as a function of the rewards, payment acceptance, and exogenous non-price characteristics.
Merchants accept the combination of cards and set retail prices to maximize their profits after paying merchant fees.
Multiproduct networks maximize profits by balancing merchant acceptance and consumer adoption.

I estimate consumer and merchant preferences by matching the reduced-form facts and several moments from payment surveys, second-choice surveys, and aggregate market-level data.
I estimate that consumers are ten times more sensitive to rewards than merchants are sensitive to fees.
A one-basis-point (1-bp) increase in Visa credit rewards increases Visa's market share among consumers by $4\%$.
In contrast, a 1-bp increase in merchant fees for Visa credit cards causes only a $0.4\%$ decline in the share of merchants that accept Visa.
This asymmetry drives high fees and rewards.

Fee caps at the cost of cash—the ``tourist test'' motivating European regulations—raise total welfare by \$29 billion annually.
Networks respond by cutting credit card rewards by 199 basis points.
With price coherence, lower merchant fees pass through to lower retail prices, benefiting low-income consumers who primarily use cash and debit.
The decline in rewards hurts high-income credit cardholders but also relieves a distortion.
In the baseline, many consumers use credit cards for the rewards despite preferring debit for its non-price characteristics.
Lower rewards allow these consumers to switch back to debit, which drives most of the welfare gain.
This \$29 billion gain dwarfs the \$12 billion welfare gain from the CARD Act \parencite{Agarwal.etal2015}, suggesting that regulating networks is at least as important as regulating issuers.

A market-based alternative is dual routing, which increases multi-homing from around 60\% to 80\% by requiring credit cards to offer merchants two processing options, as proposed in the Credit Card Competition Act. 
This change reduces credit card fees by 15 basis points and rewards by 32 basis points, increasing total welfare by \$8 billion. 

Yet without increased multi-homing, competition can worsen these problems.
I compare the status quo to a merger of all three networks into a single monopolist.
A merger to monopoly would increase total welfare by \$16 billion because the problem is not the level of competition but its direction.
Because many consumers single-home, competition worsens price coherence problems by driving adoption of high-fee, high-reward payment methods.
Revised regulation, not just more competition, is necessary to correct the market failure from price coherence.

My model also casts doubt on two other proposed solutions: caps on \emph{debit} card merchant fees and public sector entry into payment markets.
In the U.S., debit card merchant fees are capped whereas credit card fees are not.
I show that this inconsistency is regressive and reduces total welfare by pushing consumers to use credit cards.
This matches prior work on the Durbin Amendment \parencite{Mukharlyamov.Sarin2025}.
I model a public option as a debit payment network that charges no markups on merchants and pays no rewards.
The lack of rewards limits consumer adoption and thus welfare gains.

\subsection{Related Literature}

The closest related empirical work is \textcite*{Huynh.etal2022}, who estimate a structural two-sided model of consumer and merchant card adoption.
I build on their work by modeling merchant and network competition.
Merchant competition explains how credit card rewards inflate retail prices, redistribute consumption, and ultimately hurt consumers.
Network competition endogenizes merchant fees and consumer rewards, enabling me to assess how price controls and competition affect prices and total welfare.

My model combines three ingredients that prior work treats separately: consumer multi-homing, merchant heterogeneity, and merchant competition.
Assuming single-homing overstates harm from competition \parencite{Edelman.Wright2015, Anderson.etal2018, Bakos.Halaburda2020, Teh.etal2022}; merchant homogeneity rules out any positive effects of competition on merchant fees \parencite{Rochet.Tirole2011, Guthrie.Wright2007}; and ignoring merchant competition understates merchants incentives to accept cards \parencite{Rochet.Tirole2003, Wright2012,Li.etal2020}.
By combining all three features, my paper contributes to the empirical IO literature on two-sided markets \parencite{Rysman2004, Lee2013, Rosaia2020, Song2021}.

My paper also contributes to the literature on platform pricing and off-platform outcomes \parencite{Bergemann.etal2024, Chen2024}.
With price coherence, merchant fees strongly influence retail prices, explaining the welfare gains I find from fee caps.
In contrast, \textcite{Sullivan2025} finds that commission caps reduce welfare in food delivery, where price coherence is absent.

\end{document}
